\color{unidarkgreen}{\subsection{Устройство резервирования при отказе выключателя}}
\color{black}

\subsubsection{Общие сведения}

\begin{enumerate}[label=\arabic{section}.\arabic{subsection}.\arabic{subsubsection}.\arabic*,  align=left, labelsep=4pt, leftmargin=0pt, itemindent=57pt]

\item
Функциональный блок ПО УРОВ предназначен для формирования сигнала пуска УРОВ в
аварийном режиме при действии защит на отключение в цепь управления выключателем низшей
стороны трансформатора.  

\item
ПО УРОВ реагирует на превышение действующего значения основной гармоники фазного тока
выбранной стороны трансформатора или тока цепи резистора заземления нейтрали уставки \textbf{«Iср ПО
УРОВ НН1(НН)»} и на сигнал пуска от действия защит на отключение выключателя определенной
стороны (рисунок \ref{pic:UROV}). По истечении времени \textbf{«Тср ПО УРОВ НН1(НН)»} выдается сигнал срабатывания.
Предусмотрены ввод/вывод функционального блока программным переключателем \textbf{«ПО УРОВ
НН1(НН)»} и вывод с помощью внешнего сигнала.

\item
УРОВ ВН предусматривает повторное действие на себя с регулируемой выдержкой \textbf{«Тср повт
УРОВ ВН1»}, вводиться накладкой \textbf{«Действ на себя УРОВ В ВН1»}.

\item
На рисунке \ref{pic:UROV} УРОВ НН1(Н). УРОВ НН1(Н) аналогичен УРОВ НН2(СН).

\begin{figure}[H]
\centering
\includegraphics[width=0.8\textwidth,height=0.8\textheight,keepaspectratio]{\fbpath/765.0334 УРОВ/img/UROV.pdf}
\captionof{figure}{Функциональный блок УРОВ}\label{pic:UROV}
\end{figure}


\small
\begin{longtable}{|>{\centering\arraybackslash}m{5.3cm}|>{\centering\arraybackslash}m{4.2cm}|>{\centering\arraybackslash}m{1cm}|>{\centering\arraybackslash}m{3.3cm}|>{\centering\arraybackslash}m{1.8cm}|}
% Первая страница: полный заголовок
\caption{Перечень уставочных параметров УРОВ\hfill\vspace{-0.5\baselineskip}}\label{desc:tbl1}\\
\hline
\rowcolor{gray!30}
Наименование уставки & Диапазон & Шаг & Значение по умолчанию & Единица измерения \\
\hline
\endfirsthead

% Промежуточные страницы: «Продолжение таблицы»
\caption*{\hspace{3pt}\emph{Продолжение таблицы \ref{desc:tbl1}}\hfill\ }\vspace{-0.5\baselineskip} \\
\hline
\rowcolor{gray!30}
Наименование уставки & Диапазон & Шаг & Значение по умолчанию & Единица измерения \\
\hline
\endhead

%===m>ДТЗ|-
%===m
\end{longtable}
\normalsize	

\end{enumerate}